% Options for packages loaded elsewhere
\PassOptionsToPackage{unicode}{hyperref}
\PassOptionsToPackage{hyphens}{url}
\documentclass[
  ignorenonframetext,
]{beamer}
\newif\ifbibliography
\usepackage{pgfpages}
\setbeamertemplate{caption}[numbered]
\setbeamertemplate{caption label separator}{: }
\setbeamercolor{caption name}{fg=normal text.fg}
\beamertemplatenavigationsymbolsempty
% remove section numbering
\setbeamertemplate{part page}{
  \centering
  \begin{beamercolorbox}[sep=16pt,center]{part title}
    \usebeamerfont{part title}\insertpart\par
  \end{beamercolorbox}
}
\setbeamertemplate{section page}{
  \centering
  \begin{beamercolorbox}[sep=12pt,center]{section title}
    \usebeamerfont{section title}\insertsection\par
  \end{beamercolorbox}
}
\setbeamertemplate{subsection page}{
  \centering
  \begin{beamercolorbox}[sep=8pt,center]{subsection title}
    \usebeamerfont{subsection title}\insertsubsection\par
  \end{beamercolorbox}
}
% Prevent slide breaks in the middle of a paragraph
\widowpenalties 1 10000
\raggedbottom
\AtBeginPart{
  \frame{\partpage}
}
\AtBeginSection{
  \ifbibliography
  \else
    \frame{\sectionpage}
  \fi
}
\AtBeginSubsection{
  \frame{\subsectionpage}
}
\usepackage{iftex}
\ifPDFTeX
  \usepackage[T1]{fontenc}
  \usepackage[utf8]{inputenc}
  \usepackage{textcomp} % provide euro and other symbols
\else % if luatex or xetex
  \usepackage{unicode-math} % this also loads fontspec
  \defaultfontfeatures{Scale=MatchLowercase}
  \defaultfontfeatures[\rmfamily]{Ligatures=TeX,Scale=1}
\fi
\usepackage{lmodern}
\ifPDFTeX\else
  % xetex/luatex font selection
\fi
% Use upquote if available, for straight quotes in verbatim environments
\IfFileExists{upquote.sty}{\usepackage{upquote}}{}
\IfFileExists{microtype.sty}{% use microtype if available
  \usepackage[]{microtype}
  \UseMicrotypeSet[protrusion]{basicmath} % disable protrusion for tt fonts
}{}
\makeatletter
\@ifundefined{KOMAClassName}{% if non-KOMA class
  \IfFileExists{parskip.sty}{%
    \usepackage{parskip}
  }{% else
    \setlength{\parindent}{0pt}
    \setlength{\parskip}{6pt plus 2pt minus 1pt}}
}{% if KOMA class
  \KOMAoptions{parskip=half}}
\makeatother
\usepackage{longtable,booktabs,array}
\newcounter{none} % for unnumbered tables
\usepackage{calc} % for calculating minipage widths
\usepackage{caption}
% Make caption package work with longtable
\makeatletter
\def\fnum@table{\tablename~\thetable}
\makeatother
\setlength{\emergencystretch}{3em} % prevent overfull lines
\providecommand{\tightlist}{%
  \setlength{\itemsep}{0pt}\setlength{\parskip}{0pt}}
% Auto-generated by infrastructure/rendering/_pdf_latex_helpers.py::extract_math_font_preamble.
% Minimal header passed to Pandoc via -H so Beamer slides pick up the same math font
% as the combined-PDF path. Adjust the manuscript's preamble.md to change the font globally.
\usepackage{unicode-math}
\setmathfont{latinmodern-math.otf}

% Slide layout defaults for warning-clean scientific decks.
\usepackage{etoolbox}
\IfFileExists{xurl.sty}{\usepackage{xurl}}{}
\IfFileExists{seqsplit.sty}{\usepackage{seqsplit}}{\newcommand{\seqsplit}[1]{#1}}
\protected\def\breakseq#1{\seqsplit{#1}}
\protected\def\breaktt#1{\begingroup\ttfamily\seqsplit{#1}\endgroup}
\setlength{\emergencystretch}{6em}
\tolerance=5000
\hbadness=10000
\hfuzz=1pt
\setlength{\tabcolsep}{2pt}
\AtBeginEnvironment{longtable}{\tiny\renewcommand{\arraystretch}{0.86}\setlength{\tabcolsep}{1pt}}
\AtBeginEnvironment{tabular}{\tiny\renewcommand{\arraystretch}{0.86}\setlength{\tabcolsep}{1pt}}

% Natbib and cross-reference fallbacks — slides don't load natbib
% or cleveref, but combined-PDF manuscript prose may emit these
% commands. The fallback renders citations as a bracketed key list
% and cross-references as detokenized labels so slides don't fail on
% undefined control sequences or raw underscores. \providecommand is
% a no-op if packages load later.
\providecommand{\citep}[1]{[#1]}
\providecommand{\citet}[1]{#1}
\providecommand{\citealp}[1]{#1}
\providecommand{\citeauthor}[1]{#1}
\providecommand{\citeyear}[1]{#1}
\providecommand{\cref}[1]{\texttt{\detokenize{#1}}}
\providecommand{\Cref}[1]{\texttt{\detokenize{#1}}}
\usepackage{bookmark}
\IfFileExists{xurl.sty}{\usepackage{xurl}}{} % add URL line breaks if available
\urlstyle{same}
% fallback for those not using the hyperref driver hyperxmp:
\makeatletter
\@ifundefined{xmpquote}{\newcommand{\xmpquote}[1]{#1}}{}
\makeatother
\hypersetup{
  hidelinks,
  pdfcreator={LaTeX via pandoc}}

\author{\texorpdfstring{}{}}
\date{}

\begin{document}

\section{Authoring Contract: Human Review and Forking
Obligations}\label{authoring-contract-human-review-and-forking-obligations}

\begin{frame}[allowframebreaks]{Authoring Contract: Human Review and
Forking Obligations}
The authoring contract is configured to state human responsibilities,
name fork obligations, connect review to generated evidence, require
domain validators before domain claims, and document fork migration
notes. It addresses pipeline maintainers directly because the generator
can preserve structure, provenance, and reviewability, but it cannot
decide what a field should claim. Human authors remain responsible for
theory, interpretation, citations, and the choice to publish.

The declared obligations are Review generated claims, Review config
diffs, Extend claim evidence, Add domain validators, Rerun the full
project path, Review method invariants, Assemble reviewer packet, Write
fork migration notes. They convert responsible use into a checklist:
review generated claims in the hydrated manuscript, inspect the
generated evidence tables, extend the claim ledger when new claims
appear, and add domain-specific validators before the template is used
for a real empirical or theoretical manuscript.

The quality standard is claim humility. A fork that only changes words
has not preserved the exemplar. A responsible fork changes the config,
adjusts source-owned composition where necessary, regenerates artifacts,
and reruns the same tests and validation gates before treating the
output as reader-ready.
\end{frame}

\begin{frame}[fragile,allowframebreaks]{Authoring Obligations}
\protect\phantomsection\label{authoring-obligations}
{\def\LTcaptype{none} % do not increment counter
\begin{longtable}[]{@{}
  >{\raggedright\arraybackslash}p{(\linewidth - 4\tabcolsep) * \real{0.3333}}
  >{\raggedright\arraybackslash}p{(\linewidth - 4\tabcolsep) * \real{0.3333}}
  >{\raggedright\arraybackslash}p{(\linewidth - 4\tabcolsep) * \real{0.3333}}@{}}
\toprule\noalign{}
\begin{minipage}[b]{\linewidth}\raggedright
Obligation
\end{minipage} & \begin{minipage}[b]{\linewidth}\raggedright
Required action
\end{minipage} & \begin{minipage}[b]{\linewidth}\raggedright
Review surface
\end{minipage} \\
\midrule\noalign{}
\endhead
Review generated claims & Inspect hydrated manuscript bodies before
copied outputs are treated as reader-ready. &
\breaktt{output/manuscript\ and\ output/web} \\
Review config diffs & Treat lexicon, slot, title, move, and
section-switch edits as source-data changes. &
\breaktt{manuscript/config.yaml} \\
Extend claim evidence & Update the claim ledger when generated prose
adds a new claim boundary. & \breaktt{data/claim\_ledger.yaml} \\
Add domain validators & Add tests and validation artifacts before using
the template for domain-specific claims. &
\breaktt{tests\ and\ output/reports} \\
Rerun the full project path & Regenerate analysis artifacts, render
outputs, validate outputs, and copy deliverables. &
\breaktt{pipeline\ command\ logs} \\
Review method invariants & Confirm changed tokens are explained only by
seed, slot, category, ordinal, or category inventory changes. &
\breaktt{src/tokens.py,\ token\_inventory.json,\ and\ output/manuscript/02\_methodology.md} \\
Assemble reviewer packet & Provide hydrated manuscript, rendered
outputs, figures, data, reports, validation report, and output
statistics together. & \breaktt{output/templates/template\_madlib} \\
Write fork migration notes & Document config, source, test, validator,
and claim-ledger changes required by a domain fork. &
\breaktt{README.md,\ STANDALONE.md,\ and\ data/claim\_ledger.yaml} \\
\bottomrule\noalign{}
\end{longtable}
}
\end{frame}

\end{document}
